\chapter{绪论}

\section{课题来源}
本课题来源于国家自然科学基金集成项目“面向航空制造的人-机器人协作技术及应用研究”，资助号：92048301。

\section{课题研究背景及意义}

在国家航空装备智能制造升级需求推动下，机舱电气系统的自动化装配已成为提升整机质量与可靠性的重要方向。电连接器作为机舱内信号与能量传输的关键节点，其装配质量直接关系到飞行安全与维护成本。当前该工序仍以人工为主，操作者需在布满管路、线槽与设备舱壁的狭窄空间中借助手持工具完成插接与锁紧，长期高强度作业导致一致性、可重复性和过程可追溯性不足。随着机舱内部系统集成度持续提高，单纯依赖人工已难以满足高可靠、标准化装配的要求，机器人替代或协同人工作业具有迫切的工程需求。

航空机舱装配环境与传统工业产线存在显著差异，可归纳为三方面。第一，空间约束突出：机舱内部结构件、舱壁与管线支架密集，布线区域狭窄且多处于深腔位置，机械臂接近目标时面临典型窄通道问题。第二，光照条件恶劣：金属舱壁强反射、局部照明不均与深阴影并存，严重干扰基于边缘与纹理特征的传统视觉算法。第三，装配对象呈刚柔耦合特性：电连接器尾部整合有柔性线缆，线缆在重力与布线约束作用下呈随机弯曲形态，在视觉层面遮挡壳体、定位孔等关键几何特征，在动力学层面引入随位姿变化的拖拽力与回复力矩，使感知精度与装配可靠性同时受到制约。

上述环境特殊性决定了电连接器机器人自主抓取与装配本质上是一个涉及复杂环境感知、柔性障碍建模、狭窄空间规划和不确定接触控制的跨模块耦合问题。然而现有研究大多针对单一子问题展开，如仅优化6D位姿估计精度、仅考虑刚性障碍下的运动规划，或仅围绕标准peg-in-hole任务进行策略学习，尚未形成针对航空机舱环境中带柔性拖缆电连接器抓取与装配的系统化解决方案。

在此背景下，构建面向航空制造的电连接器机器人自主抓取与装配系统，具有重要的理论价值和工程意义。在理论层面，该研究有助于推动复杂遮挡下目标感知与位姿估计、柔性障碍参与的受限空间规划以及带未知扰动的接触装配控制等方向的交叉融合，为构建"感知—规划—控制"闭环的机器人装配理论体系提供实践支撑。在工程层面，通过实现电连接器从识别、抓取、避障搬运到柔顺插接的全流程自动化，可显著提升航空机舱装配效率与一致性，降低人工劳动强度，并为后续复杂线缆束装配与人机协作奠定基础。
\section{国内外研究现状}

\subsection{目标感知与 6D 位姿估计研究现状}

目标感知与 6D 位姿估计是机器人自主装配系统的前端关键环节，其任务是从单目、双目或 RGB-D 观测中获取目标类别、实例身份及完整的六自由度姿态信息，为后续抓取规划与装配控制提供几何约束。在航空机舱电连接器装配场景中，该问题因多实例共存、金属壳体纹理贫乏、局部强反光以及柔性线缆遮挡等因素而呈现更高的复杂度，不仅需要完成连接器身份识别与精确位姿估计，还需同步提取线缆延伸形态以服务于后续规划与力控模块。

早期 6D 位姿估计方法以 EPnP\cite{Lepetit2009} 等几何求解器与 RANSAC 鲁棒估计范式为代表，在已知二维—三维对应关系条件下可高效求解物体位姿，但高度依赖可靠的局部特征提取。在航空机舱环境中，金属壳体弱纹理与强镜面反射使特征检测器频繁失效，柔性线缆的大面积遮挡进一步削弱了可用特征，导致几何方法鲁棒性明显下降。

随着深度学习的快速发展，基于卷积神经网络的目标检测方法成为复杂场景视觉感知的主流方案。以 Faster R-CNN\cite{Ren2015} 为代表的两阶段检测框架（如图\ref{fig:ch1:faster-rcnn}所示）与以 YOLO\cite{Redmon2016} 为代表的单阶段回归式检测框架（如图\ref{fig:ch1:yolo}所示）分别从精度和速度两个维度推动了目标检测的进步。Wang 等人于 2023 年提出的 YOLOv7\cite{Wang2023yolov7} 进一步将实时检测精度推至新高度。

\begin{figure}[htbp]
  \centering
  \fbox{\parbox[c][0.20\textheight][c]{0.85\textwidth}{\centering 图1-1\\（空白占位图）}}
  \caption{Faster R-CNN 模型结构示意图（RPN 区域建议网络与 RoI Pooling 联合检测框架）}
  \label{fig:ch1:faster-rcnn}
\end{figure}

\begin{figure}[htbp]
  \centering
  \fbox{\parbox[c][0.20\textheight][c]{0.85\textwidth}{\centering 图1-2\\（空白占位图）}}
  \caption{YOLO 模型结构示意图（单阶段回归式检测框架的端到端预测流程）}
  \label{fig:ch1:yolo}
\end{figure}

在实例分割方面，Mask R-CNN\cite{He2017maskrcnn} 通过在检测框架中增加像素级掩膜预测分支开创了检测—分割联合学习范式；DeepLabV3+\cite{Chen2018deeplabv3plus} 则以编码器—解码器结构与空洞可分离卷积在语义分割精度与效率之间取得了良好平衡。上述技术的持续进步为复杂遮挡条件下的电连接器多实例识别与线缆掩膜提取奠定了方法基础。

在基于深度学习的 6D 位姿估计方面，研究路线已从 PoseCNN\cite{Xiang2018posecnn} 等直接回归方法发展为关键点投票、稠密对应和多模态融合等多条技术路径。Peng 等人于 2019 年提出 PVNet\cite{Peng2019pvnet}，采用像素级关键点投票策略，通过对目标表面各像素预测指向预定义关键点的方向向量并进行投票聚合来定位关键点，有效缓解了遮挡对关键点检测的影响。Wang 等人于 2019 年提出 DenseFusion\cite{Wang2019densefusion}，将 RGB 特征与深度点云特征在像素级进行稠密融合，并通过迭代精化网络逐步优化位姿预测，在 YCB-Video 等基准数据集上取得了当时的最优精度。

He 等人于 2020 年提出 PVN3D\cite{He2020pvn3d}，将关键点投票机制拓展至三维点云空间，通过深度逐点投票网络在三维空间中定位关键点，进一步提升了对遮挡和杂乱背景的鲁棒性。Song 等人于 2020 年提出 HybridPose\cite{Song2020hybridpose}，引入关键点、边向量和对称对应三种混合中间表示，利用多类型几何特征的互补性增强位姿回归在遮挡工况下的稳定性。上述方法的技术路线对比如图\ref{fig:ch1:pose-methods}所示。

\begin{figure}[htbp]
  \centering
  \fbox{\parbox[c][0.22\textheight][c]{0.85\textwidth}{\centering 图1-3\\（空白占位图）}}
  \caption{典型 6D 位姿估计方法技术路线对比示意图（由直接回归、关键点投票、稠密对应到多模态融合的演进路径）}
  \label{fig:ch1:pose-methods}
\end{figure}

在多实例场景与面向装配的位姿估计方面，近年来涌现了若干重要进展。He 等人于 2021 年提出 FFB6D\cite{He2021ffb6d}，构建了 RGB 与点云之间的双向全流融合网络，在特征编码阶段实现两种模态信息的深度交互，显著提升了对低纹理与遮挡物体的位姿估计精度。Labb\'{e} 等人于 2020 年提出 CosyPose\cite{Labbe2020cosypose}，面向多视图多物体场景，通过单视图位姿假设生成、跨视图鲁棒匹配以及物体级光束法平差三个阶段的级联处理，实现了对多实例场景中各物体 6D 位姿的一致性估计，并在 BOP Challenge 2020 中获得最佳整体方法与最佳纯 RGB 方法等多项冠军。Su 等人于 2022 年提出 ZebraPose\cite{Su2022zebrapose}，采用分层二进制表面编码策略，以由粗到精的方式建立图像像素与模型表面的稠密对应关系，在 LM-O 与 YCB-V 基准上部分指标已超越 RGB-D 方法。上述研究表明，6D 位姿估计技术正从单对象静态识别向面向多实例、多模态融合和装配任务导向的综合感知阶段演进。

在柔性线缆与细长变形物体感知方面，Tang 等人于 2018 年提出基于相干点漂移的柔性线性物体操作框架\cite{Tang2018}，利用 B 样条曲线建模线缆形状并结合点云信息实现三维状态估计，为线缆形态的实时感知提供了方法基础。Caporali 等人于 2023 年提出 DLO3DS 方法\cite{Caporali2023}，通过多视图二维图像重建柔性线性物体的三维形状，平均重建误差达到 0.82 mm，验证了在工业级精度要求下利用普通相机实现线缆三维感知的可行性。

Mou 等人于 2022 年在重型工业电缆操作研究中提出基于二维卷积神经网络的电缆耦合效应建模方法\cite{Mou2022}，通过数据驱动方式学习电缆对机器人末端力与力矩的影响规律，耦合效应估计精度超过 85\%，表明将柔性电缆的动力学特性纳入系统建模可显著提升对耦合干扰的预测能力。

在国内研究方面，清华大学 Li 等人于 2019 年提出 CDPN\cite{Li2019cdpn}，通过将旋转与平移解耦为坐标预测与位姿回归两个独立子任务，在 LINEMOD 数据集上取得了当时最优的纯 RGB 位姿估计精度，并在 BOP Challenge 2019 中获最佳纯 RGB 方法奖，验证了解耦策略在弱纹理目标位姿估计中的有效性。在工业视觉检测方面，Wu 等人于 2020 年提出基于改进 YOLOv3 的电连接器机器视觉检测方法\cite{Wu2020connector}，通过 K-means 聚类优化锚框比例、特征图融合与残差结构改进，实现了对电连接器焊点缺陷与旋转缺陷的自动化检测，准确率达 93.5\%，表明基于深度学习的目标检测方法已具备在电连接器工业检测场景中落地的可行性。上述国内工作为本文面向航空机舱场景的电连接器多实例检测与线缆精细分割提供了重要的方法借鉴与工程经验。

综合上述研究可以看出，目标感知与 6D 位姿估计技术已在遮挡鲁棒性、多模态融合和多实例处理方面取得显著进展，但针对"柔性线缆严重遮挡条件下多实例电连接器区分与面向装配的位姿及线缆形态联合估计"这一组合问题仍缺乏针对性解决方案。现有研究通常将柔性线缆视为背景或一般遮挡源，未将其系统建模为后续抓取规划与力控装配所需的重要几何与物理信息载体，这限制了感知结果在整条抓取—装配链路中的可利用程度。

\subsection{复杂受限空间规划与避障研究现状}

在完成目标感知与位姿估计之后，机器人需要在复杂受限空间内规划从当前构型到目标抓取位姿以及从抓取位姿到装配位姿的无碰撞轨迹，运动规划算法的性能直接影响系统的安全性与效率。传统路径规划方法包括基于图搜索的 A* 与 D* 算法，以及基于随机采样的 PRM 和 RRT 系列算法，其中 RRT 类方法因对高维配置空间的良好适应性而被广泛用于机械臂运动规划问题。

Kuffner 与 LaValle 于 2000 年提出 RRT-Connect 算法\cite{Kuffner2000}，采用双树同步扩展与贪心连接策略使高维配置空间中的规划效率显著提升，已成为机械臂运动规划的经典基线方法，其搜索树扩展过程如图\ref{fig:ch1:rrt-connect}所示。然而后续研究表明，该算法在存在长窄通道的配置空间中仍面临大量无效采样问题，桥接测试等采样策略虽可改善窄通道场景中的可行性，但整体收敛速度仍有提升空间。

\begin{figure}[htbp]
  \centering
  \fbox{\parbox[c][0.22\textheight][c]{0.85\textwidth}{\centering 图1-4\\（空白占位图）}}
  \caption{RRT-Connect 算法双树搜索扩展示意图（起始树 $\mathcal{T}_\mathrm{init}$ 与目标树 $\mathcal{T}_\mathrm{goal}$ 交替扩展并贪心连接的过程）}
  \label{fig:ch1:rrt-connect}
\end{figure}

在渐近最优规划方面，Karaman 与 Frazzoli 于 2011 年提出的 RRT*\cite{Karaman2011} 通过代价函数与邻域重连机制从理论上证明了路径代价的渐近最优性，Gammell 等人在此基础上先后提出 Informed RRT*\cite{Gammell2014} 与 BIT*\cite{Gammell2015}，分别通过椭球启发式采样和批量启发式搜索进一步加速了收敛。在轨迹优化方面，CHOMP\cite{Zucker2013} 利用协变量哈密顿框架在轨迹平滑性与障碍物代价间进行迭代优化，为受限空间中的路径质量优化提供了有效手段。

在面向特定受限环境的规划改进方面，已有研究针对工业装配场景中的狭窄空间约束展开工作。相关方法通过引入目标偏置采样、环境势场引导和窄通道检测等启发式策略，在典型窄通道场景中显著缩短规划时间并提高路径成功率。同时，面向多约束环境的规划方法将几何约束、关节极限和末端姿态要求统一纳入采样与校验框架，为航空制造等空间高度受限的应用场景提供了技术参考。上述研究表明，通过引入启发式信息、环境结构特征以及局部优化策略，可显著改善传统 RRT 类算法在受限空间中的性能表现。

在抓取生成方面，近年来基于深度学习的六自由度抓取生成方法取得了快速发展。Mousavian 等人于 2019 年提出 6-DOF GraspNet\cite{Mousavian2019}，采用变分自编码器从单目深度图像生成六自由度抓取姿态，通过对抓取质量进行评估与精化实现了对未知物体的泛化抓取生成。Fang 等人于 2020 年构建 GraspNet-1Billion 基准数据集\cite{Fang2020}，包含超过 97000 张 RGB-D 图像和十亿量级标注抓取位姿，并提出基于点云的端到端抓取生成网络，通过解耦接近方向与操作参数实现抓取姿态的高效预测，如图\ref{fig:ch1:grasp-gen}所示。Sundermeyer 等人于 2021 年提出 Contact-GraspNet\cite{Sundermeyer2021}，利用深度网络从场景点云直接预测接触点和六自由度抓取姿态分布，在复杂遮挡与堆叠场景中展现出良好的泛化性和较高的抓取成功率，为"感知—抓取一体化"的系统设计提供了实践范式。

\begin{figure}[htbp]
  \centering
  \fbox{\parbox[c][0.22\textheight][c]{0.85\textwidth}{\centering 图1-5\\（空白占位图）}}
  \caption{基于点云的 6DoF 抓取生成方法流程示意图（场景/局部点云输入—特征提取—接近方向与夹持参数多分支预测—候选抓取位姿集合输出）}
  \label{fig:ch1:grasp-gen}
\end{figure}

在面向装配任务的抓取规划方面，研究重点逐渐从单纯提升抓取成功率转向考虑抓取姿态与后续任务目标之间的相容性。具体而言，抓取姿态不仅需要保证夹持稳定性，还需满足后续插接操作对末端朝向和接近方向的约束，即任务相容性抓取。近期研究通过将分割模型与抓取网络级联实现对指定目标的选择性抓取，部分工作还将装配方向约束编码至抓取评分函数中，从候选集合中筛选与装配轴线对齐度最高的姿态。

在国内学者的工作中，上海交通大学卢策吾团队在通用抓取感知领域做出了系统性贡献。在构建 GraspNet-1Billion 大规模基准\cite{Fang2020}的基础上，Wang 等人于 2021 年提出基于 Graspness 先验的抓取检测方法\cite{Wang2021graspness}，通过几何可抓取性度量对场景点云进行快速筛选，在 GraspNet-1Billion 基准上将检测精度提升超过 30 AP 的同时保持了较高的推理速度。Fang 等人于 2023 年进一步提出 AnyGrasp\cite{Fang2023anygrasp}，实现了空间与时间域上鲁棒、高效的七自由度稠密抓取感知，在超过 300 种未知物体的料箱拣选任务中达到 93.3\% 的成功率，接近人类水平，为复杂杂乱场景中的高精度抓取提供了工业级解决方案。上述国内研究为本文面向航空制造场景的抓取生成与规划提供了技术基础与数据支撑。

然而，上述研究大多在开放或半结构化桌面环境中验证，鲜有工作将柔性线缆作为显式障碍纳入抓取可达性分析与路径规划。在电连接器机舱装配场景下，机器人不仅需要避开舱壁和设备等刚性障碍，还需避开与目标零件耦合的柔性线缆，其形态随时间动态变化。现有规划方法通常将柔性障碍等效为静态膨胀后的刚性体，并未充分利用视觉感知到的线缆形态信息构建动态障碍物地图，这在狭窄机舱空间内容易导致路径与线缆发生缠绕或近距离擦碰，影响抓取与装配可靠性。

\subsection{机器人精密力控装配研究现状}

机器人精密装配，尤其是 peg-in-hole、插拔连接器等接触丰富型任务，是工业机器人技术的重要应用领域之一。Whitney\cite{Whitney1982} 建立的柔顺中心理论与 Hogan\cite{Hogan1985} 提出的阻抗控制理论（其等效模型如图\ref{fig:ch1:impedance}所示）构成了力控装配的经典理论基础，在此之上衍生出导纳控制、混合力/位置控制等柔顺方法，广泛应用于处理装配过程中的接触不确定性。

\begin{figure}[htbp]
  \centering
  \fbox{\parbox[c][0.20\textheight][c]{0.85\textwidth}{\centering 图1-6\\（空白占位图）}}
  \caption{阻抗控制等效弹簧—阻尼—质量模型示意图（机器人末端与环境交互的力—位动力学关系）}
  \label{fig:ch1:impedance}
\end{figure}

随着非结构化环境下装配需求的提升，传统解析模型在应对多源不确定性与复杂干扰力时逐渐显现局限。深度强化学习因此被引入装配领域：Inoue 等人\cite{Inoue2017}率先将 DRL 应用于高精度 peg-in-hole 任务并实现亚毫米级装配，验证了数据驱动策略在精密装配中的可行性。

Vecerik 等人于 2017 年提出利用人工示教数据增强深度强化学习的探索效率\cite{Vecerik2017}，通过将示教轨迹与智能体交互数据统一存储于优先经验回放缓冲区，在稀疏奖励条件下显著加速了柔性零件插入任务的策略收敛。Beltran-Hernandez 等人于 2020 年提出基于深度强化学习的变柔顺控制方法\cite{Beltran2020}，针对位置控制型机器人的 peg-in-hole 任务通过域随机化与迁移学习实现了从仿真到实物的策略迁移，在存在孔位不确定性的条件下完成了高精度插接。

在仿真到真实迁移方面，Tobin 等人\cite{Tobin2017}提出的域随机化方法通过在仿真中随机化视觉参数实现了策略的 Sim-to-Real 迁移；Narang 等人\cite{Narang2022}提出的 Factory 仿真框架利用 GPU 加速碰撞检测将接触丰富型场景仿真速度提升超过两万倍，为大规模装配策略学习提供了高效平台。

Schoettler 等人于 2020 年针对工业连接器插入任务，提出基于深度强化学习结合视觉输入与自然奖励信号的方法\cite{Schoettler2020}，在真实机器人上完成了连接器插入操作，实验表明将强化学习与先验信息相结合可在合理的真实交互次数内求解接触丰富型插入任务。上述仿真到实机迁移的一般流程如图\ref{fig:ch1:drl-assembly}所示。

\begin{figure}[htbp]
  \centering
  \fbox{\parbox[c][0.22\textheight][c]{0.85\textwidth}{\centering 图1-7\\（空白占位图）}}
  \caption{基于深度强化学习的装配策略仿真到实机迁移流程示意图（仿真环境构建—域随机化训练—Sim-to-Real 策略迁移—真实机器人部署）}
  \label{fig:ch1:drl-assembly}
\end{figure}

在多模态策略学习与奖励设计方面，研究者致力于将视觉、力觉和本体感觉信息统一纳入装配策略的学习框架。Luo 等人于 2021 年发表面向工业装配的鲁棒多模态策略大规模研究\cite{Luo2021}，通过系统比较结合示教的 DRL 方法与专业工业集成商的基线方案，在 NIST 标准装配基准上证明 DRL 系统在速度与可靠性方面均优于后者，为 DRL 在真实工业场景中的落地提供了有力证据。Jiang 与 Chen 于 2023 年提出结合 DDPG 与模糊动作生成器的大直径轴孔装配策略\cite{Jiang2023}，通过模糊规则补偿强化学习输出的动作偏差，在保持训练收敛速度的同时将装配接触力降低约 30\%，为复杂接触工况下精密装配的奖励函数与动作空间设计提供了参考。

在柔性电缆干扰方面，Mou 等人于 2022 年提出基于学习的电缆耦合效应建模方法\cite{Mou2022}，通过二维卷积神经网络预测重型工业电缆对机器人末端的附加力与力矩，耦合效应估计精度超过 85\%，末端定位误差小于 0.01 mm，表明柔性线缆对运动控制及装配过程具有不可忽视的动力学影响，将其显式纳入控制模型可显著提升操作精度。

在国内相关研究方面，Chen 等人于 2023 年提出基于组合强化学习的机器人轴孔装配主动柔顺控制方法\cite{Chen2023active}，通过状态判断深度 Q 网络与参数优化深度 Q 网络的级联架构在线调整变阻抗控制器参数，在 0.40 mm 间隙条件下装配成功率达 96\%，验证了将强化学习策略分层解耦可有效提高物理机器人装配任务的学习效率与成功率。结合前述 Jiang 与 Chen\cite{Jiang2023}在模糊策略补偿以及 Mou 等人\cite{Mou2022}在柔性电缆耦合效应建模方面的工作，国内学者在强化学习与力控装配的交叉领域已形成多点布局，但针对带拖缆电连接器这一特殊对象的干扰补偿策略仍有待深入。

然而，现有基于 DRL 的装配策略大多将外部干扰力简化为随机噪声或统计分布，并未显式利用视觉信息估计柔性线缆的方向与形态来构建前馈补偿项。在带拖缆电连接器装配任务中，线缆拖拽力具有明确的方向性且随插接深度非线性变化，仅依靠策略网络的隐式建模难以高效应对此类结构化干扰，限制了现有方法在刚柔耦合装配场景中的适用性。

\subsection{国内外研究现状分析}

综合上述三个方向的研究进展，目标感知与位姿估计、受限空间运动规划以及精密力控装配等领域均已形成较为丰富的方法体系，但现有成果大多围绕单一子问题独立展开，针对航空机舱环境下带柔性拖缆电连接器的多模块紧密耦合工程场景，尚未形成系统性解决方案。在视觉感知方面，基于关键点投票与多模态融合的 6D 位姿估计方法在遮挡鲁棒性上持续提升，但主要面向刚性遮挡工况设计，航空机舱中柔性线缆造成的大面积、形态时变遮挡与之存在本质差异；同时线缆延伸方向与空间形态尚未被纳入感知输出，导致后续规划与力控环节缺乏完整的几何先验。在运动规划与抓取方面，RRT* 系列算法与深度学习抓取生成网络在开放或半结构化场景中表现良好，但柔性线缆构成的动态障碍难以通过传统膨胀策略准确刻画，现有抓取评分体系也较少纳入后续装配方向约束，抓取—装配链路的任务相容性缺乏系统性考量。在力控装配方面，阻抗控制与深度强化学习的融合已在标准接触丰富型任务中得到验证，但现有策略通常将外部干扰简化为各向同性噪声，未能刻画柔性线缆拖拽力沿特定方向的非线性变化，亟需引入视觉估计的线缆方向信息构建显式前馈补偿项。

针对上述问题，本文将围绕以下三个核心科学与工程问题展开研究：（1）在复杂光照、多实例和柔性线缆严重遮挡条件下，如何构建能够同时输出电连接器身份、6D 位姿以及线缆延伸方向等信息的级联视觉感知系统，为后续抓取与装配提供高质量输入。（2）在狭窄机舱空间中，如何利用视觉感知到的舱壁与线缆形态信息构建动态障碍物地图，并在此基础上设计兼顾搜索效率与安全性的引导式 RRT-Connect 规划策略，实现面向装配的连接器抓取路径规划。（3）在存在柔性拖缆显著干扰的装配过程中，如何利用视觉估计的线缆方向向量构建前馈补偿项，并与深度强化学习策略相结合，形成兼具鲁棒性和安全性的柔顺装配控制框架。

\section{本文的主要研究内容}

针对航空机舱电连接器机器人自主抓取与装配任务中柔性线缆遮挡、狭窄空间约束与刚柔耦合干扰等关键问题，本文从视觉感知、抓取规划、柔顺装配及系统集成四个方面展开研究，构建面向航空制造的电连接器机器人自主抓取与装配系统。课题主要研究内容与全文组织架构如图\ref{fig:ch1:overview}所示，具体包括以下四个方面。

\begin{figure}[htbp]
  \centering
  \includegraphics[width=0.95\textwidth]{C:/Users/31724/.cursor/projects/d/assets/c__Users_31724_AppData_Roaming_Cursor_User_workspaceStorage_ecdb843983927f0bd3607f2900fc4e87_images_image-10e2022c-d4c9-4334-8521-54d63cead3b3.png}
  \caption{课题主要研究内容流程图（级联视觉感知—自主抓取规划—柔顺装配控制—系统集成验证的级联关系与各章对应）}
  \label{fig:ch1:overview}
\end{figure}

（1）面向柔性线缆遮挡环境的电连接器级联视觉感知与 6D 位姿估计方法研究。针对多实例混叠、光照复杂及柔性线缆遮挡等问题，设计"由粗到精"的级联视觉感知框架：第一阶段利用 YOLOv8-seg 完成多电连接器的检测、标签识别与实例分割；第二阶段在裁剪后的兴趣区域上采用改进 DeepLabV3+ 进行关键点热力图预测与线缆精细分割；最后结合 EPnP/RANSAC 输出电连接器 6D 位姿及线缆方向向量，为后续模块提供统一感知接口。

（2）面向狭窄机舱空间的电连接器自主抓取与引导式 RRT-Connect 规划策略研究。以视觉感知输出的位姿与线缆掩膜为输入，构建基于点云的端到端抓取生成网络，并通过兼顾抓取质量、线缆安全间隙与装配轴线对齐度的综合评分函数选取最优抓取姿态。在路径规划方面，提出引导式 RRT-Connect 算法，将舱壁与线缆形态映射至配置空间构建混合障碍物地图，规划避开柔性线缆的安全搬运轨迹。

（3）面向柔性拖缆干扰的深度强化学习柔顺装配控制策略研究。针对线缆拖拽产生的方向性时变干扰力，构建"视觉前馈 + 导纳控制 + 深度强化学习"的双层装配控制架构：在动力学层面引入视觉估计的线缆方向向量构建前馈补偿项，在策略层面利用 DRL 算法对位置误差、接触力及线缆方向进行联合感知，学习不同干扰工况下的柔顺装配策略。

（4）电连接器机器人自主抓取与装配系统集成及实验验证。搭建包含工业机械臂、视觉传感器和力/力矩传感器的实验平台，完成各模块系统集成，分别开展视觉感知、抓取规划与柔顺装配实验，通过多工况对比测试验证所提方法的有效性与工程应用潜力。

全文共分五章：第一章为绪论，阐述课题背景、研究现状与主要研究内容；第二章研究级联视觉感知与 6D 位姿估计方法；第三章研究自主抓取生成与引导式运动规划策略；第四章研究基于深度强化学习的柔顺装配控制方法；第五章完成系统集成与综合实验验证。

